% P8 reproducible telemetry detector: independent re-derivation from raw reward tensors
\subsection{Reproducible re-derivation of the telemetry detector}
\label{sec:p8-reproducible-detector}

The earlier integrity-detector figure in this paper was reported as a
prose-level claim (full-telemetry vs.\ reward-only AUROC). This
subsection replaces that claim with an \emph{independent, end-to-end
reproducible} re-derivation: a single script
(\texttt{experiments/openings/p8\_detector.py}) that reads the raw
per-step reward tensors under
\texttt{experiments/results/n2\_reward\_tensor\_resume/} and recomputes
every feature from the underlying \texttt{rewards}, \texttt{lengths},
and \texttt{advantages} arrays -- no pre-baked scalar field is trusted.
The task is to separate genuine RL (\texttt{grpo}) from three surrogate
methods (\texttt{aero}, \texttt{areal}, \texttt{gift}); positive
$=$ surrogate. Each of the $4$ methods contributes $40$ training steps,
giving $N{=}160$ rows with a $3{:}1$ class imbalance ($120$ surrogate
positives, $40$ genuine negatives). The classifier is a
\texttt{SimpleImputer(median)}$\to$\texttt{StandardScaler}$\to$%
\texttt{LogisticRegression} pipeline scored by out-of-fold predictions
under $5$-fold stratified cross-validation.

\begin{table}[t]
\centering
\small
\caption{Reproducible integrity-detector AUROC on the $N{=}160$
reward-tensor rows ($120$ surrogate / $40$ genuine). ``full telemetry''
uses all eleven honestly-recomputed features; ``reward-only'' uses the
three reward-level features a naive auditor would inspect. AUROC is the
out-of-fold score under $5$-fold stratified CV; the seed columns
summarise stability across CV seeds $\{0,1,2,3,4\}$.}
\label{tab:p8-reproducible-detector}
\begin{tabular}{lccc}
\toprule
Feature set & AUROC (seed $0$) & AUROC (mean $\pm$ sd, $5$ seeds) & \# features \\
\midrule
Full telemetry & $0.826$ & $0.838 \pm 0.010$ & $11$ \\
Reward-only    & $0.427$ & $0.426 \pm 0.027$ & $3$ \\
\bottomrule
\end{tabular}
\end{table}

\tableref{tab:p8-reproducible-detector} confirms the paper's qualitative
claim on independently recomputed data: \textbf{full telemetry
($0.838 \pm 0.010$) roughly doubles the reward-only baseline
($0.426 \pm 0.027$)}, and the reward-only classifier is at or below
chance -- a naive auditor who inspects only reward levels cannot
distinguish surrogate from genuine training. The full feature set is
\texttt{reward\_mean}, \texttt{reward\_std}, \texttt{frac\_all\_zero},
\texttt{frac\_all\_one}, \texttt{zvf} (fraction of prompt-groups with
zero reward variance), \texttt{mean\_len}, \texttt{cv\_len},
\texttt{len\_reward\_corr}, \texttt{adv\_abs\_mean}, \texttt{adv\_std},
and \texttt{lag1\_autocorr} (lag-$1$ autocorrelation of the per-method
step-level \texttt{reward\_mean} series); the reward-only set is
\texttt{reward\_mean}, \texttt{frac\_all\_zero}, \texttt{frac\_all\_one}.

\paragraph{Honest scope.} This re-derivation is a confirmation, not a
verbatim reproduction, and four caveats bound it. (i) It uses a
\emph{richer} feature set than the paper's earlier figure: the group
-variance, length, and advantage-structure features carry most of the
lift, so the AUROC here is not directly comparable to the earlier
prose number. (ii) The training-side \texttt{loss} scalar mentioned in
the earlier text is \emph{omitted}: it is an optimizer-side quantity
that is not reconstructable from a single step's reward/length tensor,
so it is dropped to keep every feature self-derived and auditable.
(iii) The sample is small and imbalanced ($N{=}160$, $3{:}1$), and the
source tensors are from a single seed (seed $0$), so the absolute AUROC
should be read as a point estimate with the CV-seed spread as its only
uncertainty. (iv) The tensors do not carry base-model or task
provenance, so the detector cannot condition on -- or be confounded by
-- which model or task generated a step; closing that provenance gap is
exactly what the P5 stack-conditioning program addresses.

\paragraph{Reproduction.} \texttt{python3
experiments/openings/p8\_detector.py} regenerates
\tableref{tab:p8-reproducible-detector} in under a minute on CPU and
writes the full numeric summary (per-seed AUROC, feature lists, row
counts) to
\texttt{experiments/results/p8\_openings/detector\_reproduced.json}. No
new Tinker API calls were used.
